% \noindent\textbf{}
As a popular formulation of diffusion models~\citep{difussion_2015,ddpm,ddim,rezende2015variational,kingma2023understanding}, Flow Matching~\citep{fm1,fm2} addresses generative modeling by learning the velocity fields that transport the noise distribution $p_{\mathrm{noise}}=\mathcal{N}(0, I)$ to the data distribution $p_{\mathrm{data}}$. Given a noise point $x_0 \sim  p_{\mathrm{noise}}$, a data point $x_1 \sim p_{\mathrm{data}}$, and $t \sim \mathcal{U}[0, 1]$, a straight flow path can be defined by linear interpolation: $x_t = (1-t) x_0 + t x_1$, with corresponding velocity $v_t = x_1 - x_0$.
The same intermediate point $x_t$ can arise from different pairs of $(x_0, x_1)$, each yielding a distinct velocity $v_t$. 
Flow Matching thus trains a network $f_{\theta}(x_t, t)$ to estimate the marginal velocity field~\citep{fm1}, that is, the expected velocity conditioned on $x_t$: $v(x_t, t) = \mathbb{E}_{x_0, x_1} [v_t | x_t]$. 
In practice, this is equivalently optimized in a simpler way by fitting the empirical velocity $v_t$ obtained from randomly sampled pairs $(x_0, x_1)$ during training~\citep{fm1}:
\begin{equation}
\label{eq:fm_loss}
\vspace{-0.5em}
\mathcal{L}_{\mathrm{FM}} = \mathbb{E}_{t,x_0,x_1} \left[\left\| f_{\theta}(x_t, t) - v_t \right\|^2\right].
\vspace{-0.3em}
\end{equation}
In inference, new data can be sampled by solving the ordinary differential equation (ODE) using the learned flow model: $\frac{d x_t}{d t}=f_{\theta}(x_t, t)$.
This process is typically approximated by a numerical solver, such as Euler method, over discrete time steps:  
\begin{equation}
\vspace{-0.7em}
\label{eq:fm_ode}
x_{t_{i+1}} = x_{t_i} + (t_{i+1} - t_i)f_{\theta}(x_{t_i}, t_i).
\vspace{-0.8em}
\end{equation}

% \noindent\textbf{GAN for waveform generation.}